\section*{Розділ II. Тимчасова виборча комісія}
\addcontentsline{toc}{section}{Розділ II. Тимчасова виборча комісія}

\subsection*{2.1. Статус та функції ТВК}
\addcontentsline{toc}{subsection}{2.1. Статус та функції ТВК}
    2.1.1. Тимчасова виборча комісія (ТВК) є колегіальним органом, створеним виключно для організації та проведення перших виборів до органів студентського самоврядування Університету відповідно до цього Регламенту.

    2.1.2. ТВК діє на засадах незалежності, неупередженості, колегіальності та відкритості, керуючись принципами забезпечення рівних виборчих прав усіх студентів Університету.

    2.1.3. ТВК має виключно виборчі функції та не втручається в діяльність органів студентського самоврядування після їх обрання і початку роботи.

\subsection*{2.2. Склад та формування ТВК}
\addcontentsline{toc}{subsection}{2.2. Склад та формування ТВК}
    2.2.1. ТВК складається з 7 (семи) членів, які є студентами Університету та відповідають загальним вимогам до членів ОСС, визначеним Положенням про органи студентського самоврядування.

    2.2.2. Склад ТВК формується за принципом представництва:

        \begin{enumerate}[label=\alph*)]
            \item 3 (три) представники від факультетів/інститутів Університету;
            \item 2 (два) представники від студентів, що проживають у гуртожитках (студмістечко);
            \item 2 (два) загальноуніверситетські представники.
        \end{enumerate}

    2.2.3. Кандидати до складу ТВК висуваються шляхом самовисування.

    2.2.3.1. Відбір до складу ТВК здійснюється за категоріями, визначеними у пункті 2.2.2 цього Розділу, окремим голосуванням по кожній категорії з кількістю місць, що дорівнює квоті цієї категорії.

    2.2.3.2. У разі, якщо кількість кандидатів у категорії перевищує кількість місць, до складу ТВК включаються кандидати, які набрали найбільшу кількість голосів у межах відповідної категорії. У разі рівності голосів остаточне рішення приймається відкритим голосуванням Загальних зборів студентів між такими кандидатами.

    2.2.3.3. У разі, якщо кількість кандидатів у категорії є меншою за кількість місць, Загальні збори студентів оголошують додатковий прийом заявок тривалістю до 3 (трьох) календарних днів. Якщо після додаткового прийому залишаються вакантні місця, вони можуть бути заповнені кандидатами інших категорій за рішенням Загальних зборів студентів з урахуванням забезпечення мінімального базового представництва кожної категорії принаймні 1 (одним) членом, якщо це можливо за наявних кандидатур.

    2.2.4. Усі кандидати до складу ТВК затверджуються Загальними зборами студентів Університету простою більшістю голосів від присутніх студентів.

    2.2.5. Член ТВК не може одночасно бути кандидатом на виборах, які організовує ТВК, або представляти інтереси будь-якого кандидата чи політичної сили.

\subsection*{2.3. Керівництво ТВК}
\addcontentsline{toc}{subsection}{2.3. Керівництво ТВК}
    2.3.1. ТВК очолює Голова ТВК, який обирається членами ТВК зі свого складу простою більшістю голосів на першому засіданні після формування комісії.

    2.3.2. Повноваження Голови ТВК:

        \begin{enumerate}[label=\alph*)]
            \item Організація роботи ТВК та скликання її засідань;
            \item Представлення ТВК у взаємовідносинах з адміністрацією Університету, студентською спільнотою та громадськістю;
            \item Підписання рішень ТВК та інших офіційних документів;
            \item Координація виконання рішень ТВК;
            \item Забезпечення інформування студентської спільноти про діяльність ТВК.
        \end{enumerate}

    2.3.3. З числа членів ТВК також обирається Секретар ТВК, відповідальний за ведення документообігу, протоколювання засідань та зберігання документів комісії.

    2.3.4. У разі неможливості виконання Головою ТВК своїх обов'язків, його функції тимчасово виконує заступник Голови ТВК, який обирається одночасно з Головою ТВК. За відсутності заступника, функції Голови виконує Секретар ТВК.

\subsection*{2.4. Повноваження ТВК}
\addcontentsline{toc}{subsection}{2.4. Повноваження ТВК}
    2.4.1. ТВК має такі повноваження:

        \begin{enumerate}[label=\alph*)]
            \item Організація та проведення перших виборів до органів студентського самоврядування відповідно до цього Регламенту;
            \item Реєстрація кандидатів на виборні посади та контроль за дотриманням виборчого процесу;
            \item Прийняття рішень щодо технічних регламентів електронного голосування та процедур безпеки на власний розсуд для перших виборів;
            \item Визначення порядку підтвердження статусу студента для участі у виборах;
            \item Забезпечення прозорості та чесності виборчого процесу;
            \item Розгляд скарг та спорів, пов'язаних з проведенням перших виборів;
            \item Встановлення результатів виборів та їх оприлюднення;
            \item Передача повноважень новообраним органам студентського самоврядування.
        \end{enumerate}

    2.4.2. Для виконання своїх повноважень ТВК має право:

        \begin{enumerate}[label=\alph*)]
            \item Приймати рішення з питань організації та проведення виборів;
            \item Встановлювати процедури висування та реєстрації кандидатів;
            \item Визначати порядок та способи голосування;
            \item Створювати дільничні виборчі комісії для проведення голосування;
            \item Розглядати скарги та заяви щодо порушень виборчого процесу;
            \item Взаємодіяти з адміністрацією Університету з питань організаційного та технічного забезпечення виборів;
            \item Залучати студентів до проведення виборів на громадських засадах.
        \end{enumerate}

    2.4.3. ТВК зобов'язана:

        \begin{enumerate}[label=\alph*)]
            \item Забезпечувати рівні права та можливості для всіх кандидатів;
            \item Гарантувати прозорість та відкритість виборчого процесу;
            \item Своєчасно інформувати студентську спільноту про всі етапи виборів;
            \item Вести повну та достовірну документацію виборчого процесу;
            \item Розглядати всі скарги та звернення щодо виборів протягом встановлених строків.
        \end{enumerate}

\subsection*{2.5. Порядок роботи ТВК}
\addcontentsline{toc}{subsection}{2.5. Порядок роботи ТВК}
    2.5.1. ТВК приймає рішення на засіданнях, що проводяться за потребою під час виборчого процесу.

    2.5.2. Засідання ТВК є правомочним за присутності не менше 4 (чотирьох) членів комісії.

    2.5.3. Рішення ТВК приймаються простою більшістю голосів від присутніх членів комісії. У разі рівності голосів, голос Голови ТВК є вирішальним.

    2.5.4. Засідання ТВК протоколюються. Протокол підписується Головою та Секретарем ТВК і оприлюднюється на офіційних інформаційних ресурсах не пізніше наступного дня після засідання.

    2.5.5. ТВК може створювати робочі групи зі свого складу або із залученням інших студентів для підготовки окремих питань або проведення конкретних виборів.

\subsection*{2.6. Припинення повноважень ТВК}
\addcontentsline{toc}{subsection}{2.6. Припинення повноважень ТВК}
    2.6.1. Повноваження ТВК припиняються з моменту формування першого складу Центральної виборчої комісії студентів (ЦВКс) згідно з Положенням про Центральну виборчу комісію студентів.

    2.6.2. Про припинення своїх повноважень ТВК повідомляє студентську спільноту не пізніше ніж за 3 дні до передачі повноважень ЦВКс.

    2.6.3. ТВК передає ЦВКс всі документи, матеріали та технічні засоби, пов'язані з організацією виборів, а також звіт про проведену роботу.

    2.6.4. Повноваження окремого члена ТВК можуть бути припинені достроково у випадках:

        \begin{enumerate}[label=\alph*)]
            \item Подання особистої заяви про складення повноважень;
            \item Втрати статусу студента Університету;
            \item Систематичного невиконання обов'язків члена ТВК (за рішенням ТВК більшістю у 2/3 голосів);
            \item Порушення принципів неупередженості або конфлікту інтересів.
        \end{enumerate}

    2.6.5. У разі дострокового припинення повноважень члена ТВК, його заміщення здійснюється у порядку, визначеному для первинного формування складу ТВК.